\clearpage
\appendix
\section{附录}

\subsection{分赛区后续性能与 WNS 演进}

本节完整保留《分赛区决赛后续优化说明报告》的两套量化结果。\texttt{withMext-v2} 用于贯穿分赛区
主报告版本至后续版本的统一 RTL 仿真复测；西北赛区样例用于观察不同访存与分支组成下的效果。
两套样例的绝对时间不可横向比较，折线只展示各自内部的性能与 WNS 变化。

\begin{table}[H]
  \centering
  \caption{统一采用 withMext-v2 的完整 RTL 仿真性能演进}
  \scriptsize
  \setlength{\tabcolsep}{4pt}
  \resizebox{\textwidth}{!}{%
  \begin{tabular}{@{}lcccc@{}}
      \toprule
      优化里程碑 & 频率 & 仿真时间 & 较前基线提升 & 时序状态 \\
      \midrule
      主报告版本 & 280 MHz & 2.020 s & --- & WNS $+0.050$ ns \\
      乘法与方向预测优化 & 280 MHz & 1.943 s & $3.8\%$ & WNS $+0.031$ ns \\
      D-cache 整字写分配 & 280 MHz & 1.868 s & $3.9\%$ & WNS $+0.019$ ns \\
      load-use 晚数据旁路 & 280 MHz & 1.643 s & $12.0\%$ & WNS $-0.084$ ns \\
      迭代 RTL 除法器 & 280 MHz & 1.606910 s & $2.2\%$ & 通过 \\
      LSU 地址通路与 280 MHz 收敛 & 280 MHz & 1.606910 s & $0.0\%$ & WNS $0.000$ ns \\
      窄加载与 D-cache shadow & 280 MHz & 1.607404 s & $-0.031\%$ & WNS $0.000$ ns \\
      32 项 BTB 与更新 shadow & 280 MHz & 1.607403 s & $+0.000055\%$ & WNS $-0.011$ ns \\
      精确 ret 识别与预测 & 280 MHz & 1.607403 s & $+0.000013\%$ & WNS $-0.091$ ns \\
      部分写命中保留 & 280 MHz & 1.607403 s & $+0.000001\%$ & WNS $-0.072$ ns \\
      接口缩减与除法修正分拍 & 280 MHz & 1.607403 s & $-0.000003\%$ & WNS $-0.004$ ns \\
      \textbf{最终时序闭合结果} & \textbf{280 MHz} & \textbf{1.607403 s} & $0.0\%$ & WNS $+0.013$ ns \\
      \bottomrule
  \end{tabular}}
\end{table}

最终结果完整运行 450,072,868 周期、380,344,384 条指令，CPI 为 1.183330；相对 2.020\,s
基线缩短约 20.4\%。后期若干节点主要用于时序与资源整理，因此保留精确百分比，避免将时序优化误写
成性能收益。

\begin{table}[H]
  \centering
  \caption{西北赛区样例的性能与时序演进}
  \small
  \setlength{\tabcolsep}{5pt}
  \begin{tabular}{@{}lccc@{}}
    \toprule
    优化里程碑 & 评估时间 & 较前基线提升 & 时序状态 \\
    \midrule
    多周期单元与 280 MHz 收敛 & 6.866 s & --- & WNS $0.000$ ns \\
    窄加载与 D-cache shadow & 6.553 s & $4.6\%$ & WNS $0.000$ ns \\
    32 项 BTB 与更新 shadow & 6.508 s & $0.7\%$ & WNS $-0.011$ ns \\
    精确 ret 识别与预测 & 6.481 s & $0.4\%$ & WNS $-0.091$ ns \\
    部分写命中保留 & 6.182 s & $4.6\%$ & WNS $-0.072$ ns \\
    \textbf{接口整理与最终时序闭合} & \textbf{6.184 s} & $-0.03\%$ & WNS $+0.013$ ns \\
    \bottomrule
  \end{tabular}
\end{table}

最终结果相对 6.866\,s 基线缩短约 9.9\%。6.182\,s 到 6.184\,s 的变化来自除法修正分拍和接口
整理；每次普通除法或求余增加一拍，完整前缀中 12 次相关操作精确增加 12 周期，同时 WNS 收敛至
$+0.013$\,ns。

\begin{figure}[H]
  \centering
  \begin{tikzpicture}
    \begin{axis}[
      width=0.835\textwidth,height=6.2cm,
      symbolic x coords={主报告,乘分支,DCache,晚旁路,LSU地址通路,Shadow,BTB,Ret,部分写,最终},
      xtick=data,x tick label style={rotate=32,anchor=east,font=\scriptsize},
      ylabel={评估时间（s）},ymin=1.55,ymax=2.05,
      axis y line*=left,axis x line*=bottom,grid=major,
      legend style={at={(0,1.02)},anchor=south west,font=\scriptsize}]
      \addplot[blue!75!black,mark=*] coordinates {
        (主报告,2.020) (乘分支,1.943) (DCache,1.868) (晚旁路,1.643)
        (LSU地址通路,1.606910) (Shadow,1.607404) (BTB,1.607403) (Ret,1.607403)
        (部分写,1.607403) (最终,1.607403)};
      \addlegendentry{性能}
    \end{axis}
    \begin{axis}[
      width=0.835\textwidth,height=6.2cm,
      symbolic x coords={主报告,乘分支,DCache,晚旁路,LSU地址通路,Shadow,BTB,Ret,部分写,最终},
      xtick=\empty,ylabel={WNS（ns）},ymin=-0.11,ymax=0.07,
      axis y line*=right,axis x line=none,
      legend style={at={(1,1.02)},anchor=south east,font=\scriptsize}]
      \draw[gray,dashed] (rel axis cs:0,0.6111) -- (rel axis cs:1,0.6111);
      \addplot[red!75!black,dashed,mark=square*] coordinates {
        (主报告,0.050) (乘分支,0.031) (DCache,0.019) (晚旁路,-0.084)
        (LSU地址通路,0.000) (Shadow,0.000) (BTB,-0.011) (Ret,-0.091)
        (部分写,-0.072) (最终,0.013)};
      \addlegendentry{WNS}
    \end{axis}
  \end{tikzpicture}
  \caption{withMext-v2 的性能与 WNS 走向}
\end{figure}

\begin{figure}[H]
  \centering
  \begin{tikzpicture}
    \begin{axis}[
      width=0.835\textwidth,height=6.0cm,
      symbolic x coords={时序基线,窄加载,BTB,Ret,部分写,最终},
      xtick=data,x tick label style={rotate=25,anchor=east,font=\scriptsize},
      ylabel={评估时间（s）},ymin=6.05,ymax=6.95,
      axis y line*=left,axis x line*=bottom,grid=major,
      legend style={at={(0,1.02)},anchor=south west,font=\scriptsize}]
      \addplot[blue!75!black,mark=*] coordinates {
        (时序基线,6.866) (窄加载,6.553) (BTB,6.508) (Ret,6.481) (部分写,6.182) (最终,6.184)};
      \addlegendentry{性能}
    \end{axis}
    \begin{axis}[
      width=0.835\textwidth,height=6.0cm,
      symbolic x coords={时序基线,窄加载,BTB,Ret,部分写,最终},
      xtick=\empty,ylabel={WNS（ns）},ymin=-0.11,ymax=0.04,
      axis y line*=right,axis x line=none,
      legend style={at={(1,1.02)},anchor=south east,font=\scriptsize}]
      \draw[gray,dashed] (rel axis cs:0,0.7333) -- (rel axis cs:1,0.7333);
      \addplot[red!75!black,dashed,mark=square*] coordinates {
        (时序基线,0.000) (窄加载,0.000) (BTB,-0.011) (Ret,-0.091)
        (部分写,-0.072) (最终,0.013)};
      \addlegendentry{WNS}
    \end{axis}
  \end{tikzpicture}
  \caption{西北赛区样例的性能与 WNS 走向}
\end{figure}


\subsection{硬件加速构建配置摘要}\label{app:accel-build}

项目借助 GCC 目标选项和相关接口，按处理器提供的硬件能力配置程序镜像。启用与关闭版本采用相同输入、
迭代规模和结果判定；该配置不改变处理器 RTL。

最终构建以固定工具链、选项和程序镜像完成回归。NPC 的 10/100 次运行均通过差分与 CRC 校验，
NEMU 10,000 次运行通过相同 CRC；同时检查程序镜像摘要和构建日志，确保复测使用的二进制与
记录一致。正文据此报告硬件功能、流水边界、性能结果和验证范围。
\subsection{自动化验证截图（补充证据）}

下列两图保留分赛区阶段的可视化记录，用于展示测试输出形式。它们不替代正文中按提交标识、
任务链接、日志摘录和计数器构成的最新可复现证据。

\begin{figure}[H]
  \centering
  \includegraphics[height=0.58\textheight,trim=0 0 669bp 0,clip]{verify/image13.png}%
  \includegraphics[height=0.58\textheight,trim=0 0 644bp 0,clip]{verify/image14.png}%
  \includegraphics[height=0.58\textheight,trim=0 0 654bp 0,clip]{verify/image15.png}
  \caption{RV32I 基本指令测试结果}
\end{figure}

\begin{figure}[H]
  \centering
  \includegraphics[width=0.82\linewidth,trim=0 365bp 0 0,clip]{verify/image16.png}
  \caption{src0 仿真日志与性能统计}
\end{figure}

\subsection{测评证据索引}

\begin{table}[H]\centering
\caption{关键测评证据索引}
\scriptsize
\setlength{\tabcolsep}{3pt}
\resizebox{\linewidth}{!}{\begin{tabular}{@{}lll@{}}
\toprule
\textbf{编号}&\textbf{证据类型}&\textbf{用途}\\\midrule
E-FINAL-240&最终合入 RTL 的 240 MHz 实现归档&routed WNS $+0.014$ ns、TNS 0、无失败端点\\
E-FINAL-300&300 MHz 实现归档&routed/post-physopt WNS $-0.590/-0.580$ ns\\
E-FINAL-CUR&最终合入 RTL 实现归档&routed WNS $-0.699$ ns；最终构建配置未改变 RTL\\
E-RF-01--04&分赛区后续报告与统计数据&withMext-v2、西北样例的性能与 WNS 演进\\
E-BOARD-01&最佳阶段 UART 板测归档&6.696371860\,s、四项 CRC、板卡与 bitstream 身份\\
E-BOARD-02&正确 IP 板级复测归档&6.714950940\,s、四项 CRC、300\,MHz 实现与 bitstream 身份\\
E-BOARD-03&加速接口构建的同硬件对照&6.691154040\,s、三次一致、四项 CRC 与程序镜像身份\\
E-BOARD-COMB&组合加速板测归档&5.346451040\,s、三次一致、四项 CRC 与修正后的 COE 身份\\
E-CI-01&GitHub Actions 可复现回归&run 31680203475、提交 75329078；无硬件加速 CoreMark 5,000 次迭代及其他负载日志与计数器\\
E-CI-ACC&较早硬件加速自动化回归&run 31708291526、任务 94474591976、提交 fa30e65e；5,000 次迭代日志与完整计数器\\
E-CAND-01&当前文档对应版本&最终合入提交 8591368；包含 RTL 结构与加速接口配置\\
E-ACCEL-FINAL&最终加速构建回归&NPC 10/100 次、5.292462350\,s 拟合值、NEMU 10,000 次及镜像身份检查\\
E-PROF-CRC&CRC 热点画像与板测记录&5,840,008 次动态操作、内联溢出访存和同 ISA 板测对照\\
E-PROF-COMPUTE&定宽计算与顺序访存画像&29.16M 次位域乘法及 12.96M 个处理元素\\
E-PROF-STREAM&流式状态定向周期记录&10,000-call 软件/硬件周期及处理字节计数\\
E-PROF-DEPEND&早期依赖访存原型评估&61.2M 个读写步骤及串行访存净周期代价\\
E-PROF-BACKEND&后端分区实现审计&routed WNS/TNS、失败端点数及 16M 指令周期对照\\
E-TEST&自动化测试记录&基础指令、CPU tests 与体系结构测试输出\\
E-RTTNANO&系统移植与运行记录&Nano 3.1.5 固定版本；NEMU 命令通过、最终 RTL 启动及 D-cache 断言边界\\
\bottomrule
\end{tabular}}
\end{table}

文档外的工程归档保存精确版本标识、程序镜像摘要、bitstream 摘要和原始日志，便于团队内部复核；
正文仅保留评审理解设计与测评口径所需的信息。机器可读证据清单位于 \code{src/data/evidence.csv}。

\subsection{工程目录结构}

\makeatletter
\newcommand{\mydotfill}{%
  \leavevmode
  \leaders\hbox to 1em{\hss\raisebox{0.3ex}{\textcolor{gray}{.}}\hss}\hfill\kern0pt}
\makeatother

\noindent\begin{minipage}{0.96\textwidth}
\small
\dirtree{%
.1 jyd/ \mydotfill 工程根目录.
.2 npc/ \mydotfill CPU、SoC 封装与仿真工程.
.3 design/src/ \mydotfill 流水线、预测器、缓存、扩展与设备描述.
.3 build/gened/verilog/ \mydotfill 生成的 Verilog/SystemVerilog.
.3 sim/ \mydotfill Verilator 仿真、差分与性能统计.
.3 scripts/ \mydotfill 打包、实现、报告提取与板测脚本.
.2 nemu/ \mydotfill RISC-V 参考执行环境.
.2 sdb/ \mydotfill 差分状态、调试器与指令环.
.2 abstract-machine/ \mydotfill 裸机运行环境与 riscv32-jyd 平台.
.2 jyd-tests/rtthread-nano/ \mydotfill RT-Thread Nano 3.1.5 的 AM 移植与系统测试.
.2 jyd-vivado-proj/ \mydotfill 数字孪生 Vivado 工程、IP 与约束.
.2 doc/Finals/final-rpt/ \mydotfill 总决赛技术文档.
.3 src/chapters/ \mydotfill 分章节正文.
.3 src/tikz/ \mydotfill 独立管理的 TikZ/pgfplots 图源.
.3 src/assets/ \mydotfill 仿真截图、波形、照片与外部图片.
.3 src/data/ \mydotfill 性能、时序和证据索引数据.
}
\end{minipage}

\clearpage
\subsection{参考资料}

\small
\begin{enumerate}[label={\textup{[\arabic*]}},leftmargin=3em,itemsep=0.08em,topsep=0.2em]
  \item 全国大学生集成电路创新创业大赛组委会：《第十届集创赛作品提交要求》。
  \item 竞业达企业命题组：《竞业达企业命题作品材料参考模板》。
  \item 竞业达企业命题组：《FPGA 开发板远程连接操作手册》，2026。
  \item RISC-V International, \emph{The RISC-V Instruction Set Manual, Volume I: Unprivileged ISA}.
  \item RISC-V International, \emph{The RISC-V Instruction Set Manual, Volume II: Privileged Architecture}.
  \item AMD Xilinx, \emph{Vivado Design Suite User Guide: Design Analysis and Closure Techniques}.
  \item AMD Xilinx, \emph{7 Series FPGAs Memory Resources User Guide}.
  \item EEMBC, \emph{CoreMark: An EEMBC Benchmark},
  \href{https://github.com/eembc/coremark}{official source and run rules}.
  \item Reinhold P. Weicker, ``DHRYSTONE: A Synthetic Systems Programming Benchmark,''
  \emph{Communications of the ACM}, 27(10), 1984,
  DOI: \href{https://doi.org/10.1145/358274.358283}{10.1145/358274.358283}.
  \item Verilator Project, \emph{Verilator User's Guide},
  \href{https://verilator.org/guide/latest/}{latest online edition}.
  \item Nanjing University ProjectN, \emph{NEMU: Nanjing University Emulator},
  \href{https://github.com/NJU-ProjectN/nemu}{official repository}.
  \item RT-Thread Project, \emph{RT-Thread Nano v3.1.5},
  \href{https://github.com/RT-Thread/rtthread-nano/tree/v3.1.5}{official source archive}.
  \item RISC-V International, \emph{RISC-V Architecture Test},
  \href{https://github.com/riscv-non-isa/riscv-arch-test}{official repository}; RISC-V Software,
  \href{https://github.com/riscv-software-src/riscv-tests}{riscv-tests}.
  \item Nanjing University ProjectN, \emph{AbstractMachine},
  \href{https://github.com/NJU-ProjectN/abstract-machine}{official repository}.
  \item Nanjing University ProjectN, \emph{PA: RT-Thread 选做讲义},
  \href{https://ysyx.oscc.cc/docs/ics-pa/4.1.html\#rt-thread-\%E9\%80\%89\%E5\%81\%9A}{online course notes}.
  \item Nanjing University ProjectN, \emph{rt-thread-am},
  \href{https://github.com/NJU-ProjectN/rt-thread-am}{teaching branch adapted to AbstractMachine}.
  \item Nanjing University ProjectN, \emph{riscv-tests-am},
  \href{https://github.com/NJU-ProjectN/riscv-tests-am}{riscv-tests port for AbstractMachine}.
  \item Nanjing University ProjectN, \emph{riscv-arch-test-am},
  \href{https://github.com/NJU-ProjectN/riscv-arch-test-am}{architecture-test port for AbstractMachine}.
  \item Nanjing University ProjectN, \emph{am-kernels},
  \href{https://github.com/NJU-ProjectN/am-kernels}{AbstractMachine kernels and cpu-tests}.
  \item Chisel Project, \emph{Chisel Documentation},
  \href{https://www.chisel-lang.org/docs}{official documentation}.
  \item GitHub, \emph{Using workflow run logs},
  \href{https://docs.github.com/en/actions/how-tos/monitor-workflows/use-workflow-run-logs}{GitHub Actions documentation};
  \href{https://github.com/HanPiM/jyd/actions/runs/31680203475}{run 31680203475}.
\end{enumerate}
\normalsize

\clearpage
\subsection{AI 工具使用声明}

项目研发和技术文档整理过程中使用 OpenAI 的 ChatGPT/Codex 系列工具提供辅助。分赛区阶段主要使用
ChatGPT（GPT-4o），总决赛阶段主要使用 Codex（GPT-5 系列，Codex CLI 0.147.0）；版本号按实际使用
阶段记录，云端模型的小版本由服务端滚动更新。团队对体系结构决策、代码合入、测试结果和最终文字负责，
AI 工具不替代人工设计评审、功能验证、FPGA 实现或板级测评。

按各类工作的投入量粗略估算，AI 辅助约占：体系结构分析与优化决策 15\%，RTL/工具链/脚本代码编写与
修改 30\%，验证方案与日志整理 30\%，技术文档文字、LaTeX 排版和 TikZ 图表 60\%；折合项目总体工作量
约 25\%。这些比例表示工具参与草拟、检索或整理的工作量，不表示内容未经人工复核。主要用途如下：

\begin{itemize}
  \item \textbf{工程检索与影响分析：}辅助定位模块连接、接口定义、测试入口和实现报告中的关键路径；
  \item \textbf{脚本与重复性工作：}辅助编写 COE 转换、日志整理、报告字段提取、批量回归和文档审计脚本；
  \item \textbf{代码初稿与局部重构：}依据团队给出的结构方案生成局部实现草稿，再由团队审阅、修改和测试；
  \item \textbf{验证用例补充：}辅助整理边界条件、RAW 相关、握手反压、缓存一致性和硬件加速单元参考语义；
  \item \textbf{时序报告解读：}辅助汇总逻辑级数、布线占比、端点族和资源变化，优化方向由团队确定；
  \item \textbf{技术文档编排：}辅助章节结构、表格排版、术语一致性检查、TikZ/pgfplots 图表草拟与文字润色；
  \item \textbf{使用边界：}性能与 WNS 数据均来自仿真、Vivado 报告或板测原始记录；指令语义、接口时序和
  硬件结构均经过团队审核，不使用 AI 生成的虚构数据或未经验证的结论。
\end{itemize}
